\section{The annular pole and the two-wave estimate}
\label{sec:ex-annulus-two-wave}

The selected one-step construction needs two facts about one finite-volume
response. Its logarithmic coefficient is an integral of the response over
four-dimensional momentum space, and that response is close to the continuum
quadratic pole on an annulus. This section gives the exact polar calculation
and expands the finite-volume two-channel estimate that supplies the two
radial bounds. The remaining local chart and weighted-face estimates are
stated at their precise input interface and cross-referenced to the units
that produce them.

\subsection{From the finite response to a momentum integral}

Let $E_4$ be the fundamental momentum cube in Euclidean coordinates, with
normalized measure $c_{\mathrm{wt}}\,dx$ and
$c_{\mathrm{wt}}=(2\pi)^{-4}$. It has finite volume and contains the unit
ball. For a quadratic response picture $F$ and color coordinate $a$, let
$\mathcal K_a(x)$ be one half of the second derivative in the distinguished
response direction, with the other momentum equal to $x$, and set
$g_a(x)=\operatorname{Re}\mathcal K_a(x)$. The coefficient functional
$\alpha_a(F)$ is the second Taylor coefficient of the analytic stock
picture associated with $F$. The Fourier stock formula and differentiation
under the integral give
\[
 \alpha_a(F)=c_{\mathrm{wt}}\int_{E_4}g_a(x)\,dx.
 \tag{AT1}
\]
The common analytic bound on the differentiated symbol supplies the
integrable majorant. Equation (AT1) is an identity for the analytically
assembled response picture. The equality to a finite sampled pole-test
coefficient is a separate recording statement and is not asserted by
(AT1).

For $x\in\mathbb R^4\setminus\{0\}$ let $r=|x|$ and define the gauge pole
\[
 \Pi_G(x)=4C_G\left(
   \frac{x_1^2}{r^6}-4\frac{x_1^2x_0^2}{r^8}-\frac1{r^4}
 \right),
 \tag{AT2}
\]
where $C_G$ is the quadratic Casimir in the chosen representation. The
spherical area of $S^3$ is $2\pi^2$. Rotational invariance gives
\[
 \int_{S^3}\theta_1^2\,d\theta=\frac{2\pi^2}{4},\qquad
 \int_{S^3}\theta_1^2\theta_0^2\,d\theta=\frac{2\pi^2}{24}.
 \tag{AT3}
\]
For the second identity, set $A=\int\theta_0^4$ and
$B=\int\theta_0^2\theta_1^2$. Coordinate permutations make these
independent of the distinct indices. From $\sum_i\theta_i^2=1$,
$4A+12B=2\pi^2$. Rotating the $0$-axis to
$(e_0+e_1)/\sqrt2$ gives $A=(2A+6B)/4$, hence $A=3B$ and
$B=2\pi^2/24$. The first identity follows by integrating
$\sum_i\theta_i^2=1$.

Consequently the angular mean of the parenthesis in (AT2) is
\[
 r^{-4}\left(\frac14-\frac4{24}-1\right)
       =-\frac{11}{12}r^{-4}.
\]
Since the Euclidean polar Jacobian is $r^3dr\,d\theta$,
\[
 \int_{\mu<|x|\leq1}|x|^{-4}\,dx
       =2\pi^2\int_\mu^1\frac{dr}{r}
       =2\pi^2\log\mu^{-1}.
\]
Using $c_{\mathrm{wt}}=(2\pi)^{-4}$, we obtain the exact pole law
\[
 c_{\mathrm{wt}}\int_{\mu<|x|\leq1}\Pi_G(x)\,dx
 =-\frac{11C_G}{24\pi^2}\log\mu^{-1}
 =-\lambda_G\log\mu^{-1},
 \qquad \lambda_G=\frac{11C_G}{24\pi^2}.
 \tag{AT4}
\]
This fixes the sign and normalization of the logarithm; in particular the
pole is not defined only up to an angular constant.

\subsection{The finite-volume two-channel calculation}

We now state the two scale receipts used by the finite-volume construction.
Let $x$ be a centered momentum, $r=|x|$, and $0<\mu\leq1$. Write
$\mathcal K(x)\in\mathbb C$ for the normalized response coefficient and
$g(x)=\operatorname{Re}\mathcal K(x)$. The six complex second-derivative
traces $D,Y,c,d,c_0,d_0$ are defined so that $D$ is the interacting upper
trace, $Y$ the reference upper trace, $c$ the reference lower trace, $d$
the interacting lower trace, and $c_0,d_0$ their free halo traces. The
exact identities are
\[
 \mathcal K=\frac12\left(\frac12(D-d)-(Y-c)\right),\qquad
 \Pi_G=\frac12\left(c_0-\frac12d_0\right),
 \tag{AT5a}
\]
\[
 \mathcal K-\Pi_G=\frac12\left(\frac12D-Y\right)
       +\frac12\left((c-c_0)-\frac12(d-d_0)\right).
 \tag{AT5b}
\]
The first identity is the second derivative of the split-field response
after subtracting its reference film. The second is the free halo
identity, with marked direction and color fixed. The third is their
algebraic difference; the real-valued pole $\Pi_G$ is embedded in $\mathbb C$
in these formulas. All are exact at each finite volume.

For clarity, the second identity follows directly from the two free
propagator numerators. In the selected unit color direction $e_b$, use the
normalization $\operatorname{tr}(\operatorname{ad}(e_b)^2)=-C_G$. Put
$\rho=|x|^2$, $c_{\rm col}=-C_G$, and $L=-4x_1^2$. Along the marked
displacement $t e_0$, the common denominator is
$\Delta(t)=\rho+2x_0t+t^2$. The ghost numerator is the scalar orbital
factor $L$ times the color trace, giving $Lc_{\rm col}/\rho$. For the
vector numerator, write $p=-2i x_1$ and let $A(t)$ have entries
$A_{01}(t)=-2it$, $A_{10}(t)=2it$, and all other entries zero. Then
$p^2=L$ and direct multiplication gives
\[
 \operatorname{tr}\bigl[(pI+A(-t))(pI+A(t))\bigr]
       =4p^2-8t^2=4L-8t^2.
\]
Multiplying by the same color trace and scalar denominator factor gives
the vector numerator $\rho^{-1}(4L-8t^2)c_{\rm col}$. Differentiating
the quotients twice at $t=0$ gives
\[
 \left(\frac{Lc_{\rm col}/\rho}{\Delta}\right)''_{t=0}
       =-\frac{2Lc_{\rm col}}{\rho^3}
           +\frac{8Lc_{\rm col}x_0^2}{\rho^4},\qquad
 \left(\frac{\rho^{-1}(4L-8t^2)c_{\rm col}}{\Delta}\right)''_{t=0}
       =-\frac{16c_{\rm col}}{\rho^2}
         -\frac{8Lc_{\rm col}}{\rho^3}
         +\frac{32Lc_{\rm col}x_0^2}{\rho^4}.
\]
Thus the orbital trace minus one half of the vector trace is
\[
 \frac{8c_{\rm col}}{\rho^2}
       +\frac{2Lc_{\rm col}}{\rho^3}
       -\frac{8Lc_{\rm col}x_0^2}{\rho^4}
 =8C_G\left(\frac{x_1^2}{r^6}
       -4\frac{x_1^2x_0^2}{r^8}-\frac1{r^4}\right)
 =2\Pi_G(x),
\]
which proves the halo identity in (AT5a). The first identity is the
definition of the marked split-field second response and its reference
subtraction; (AT5b) is then obtained by subtracting the just-computed free
identity and collecting terms.

The first family of estimates is the all-centered-record receipt, with
$T_1=r+\mu$ and $R_0=1+2\pi$:
\[
 0\leq r,\quad T_1=r+\mu\leq R_0,\qquad
 |\mathcal K|\leq A_1T_1^{-4}.
 \tag{AT6}
\]
The second is the shell receipt, available only on $\mu\leq r\leq1$ and
using $T_2=r$:
\[
 |D|\leq u r^{-2}+v\mu^2r^{-6},\qquad
 |Y|\leq u_{\mathrm{ref}}r^{-2}+v_{\mathrm{ref}}\mu^2r^{-6},
 \tag{AT7}
\]
\[
 |c-c_0|\leq A_0(r+\mu^2r^{-2})r^{-4},\qquad
 |d-d_0|\leq B_0(r+\mu^2r^{-2})r^{-4}.
 \tag{AT8}
\]
All constants $A_1,u,v,u_{\mathrm{ref}},v_{\mathrm{ref}},A_0,B_0$ are nonnegative and uniform in
the finite volume and selected record. These two receipts are kept
separate: the all-centered estimate uses scale $r+\mu$, while the shell
comparison uses scale $r$. Their source producer is the common finite
family of local stationary records; its exact proof obligations are
listed in the correspondence file.

\begin{proposition}[Two wave bounds from the finite response]
\label{prop:ex-two-wave}
Suppose (AT5a)--(AT8) hold uniformly. Then
\[
 |g(x)|\leq \frac{A_1}{(r+\mu)^4},\qquad
 |g(x)-\Pi_G(x)|\leq \frac{A_2}{r^3}
                     +\frac{A_2\mu^2}{r^6}
 \quad(\mu\leq r\leq1),
 \tag{AT9}
\]
where
\[
 A_2=\max\{u+u_{\mathrm{ref}}+A_0+B_0,
                    v+v_{\mathrm{ref}}+A_0+B_0\}.
\]
\end{proposition}
\begin{proof}
The first estimate follows from (AT6) and $|\operatorname{Re}z|\leq|z|$.
On the annulus, use (AT5b) and (AT7)--(AT8):
\[
 |\mathcal K-\Pi_G|
 \leq (u+u_{\mathrm{ref}})r^{-2}
      +(v+v_{\mathrm{ref}})\mu^2r^{-6}
 +(A_0+B_0)(r^{-3}+\mu^2r^{-6}).
 \tag{AT10}
\]
Here $\tfrac12$ factors are bounded by one, and
$(r+\mu^2r^{-2})r^{-4}=r^{-3}+\mu^2r^{-6}$. Since $r\leq1$,
$r^{-2}\leq r^{-3}$. The displayed $A_2$ bounds (AT10) by
$A_2(r^{-3}+\mu^2r^{-6})$. Taking real parts gives
$|g-\Pi_G|\leq|\mathcal K-\Pi_G|$, proving (AT9).
\end{proof}

\subsection{The two applications of the finite-volume receipts}

The common family of local records is used twice, with different scale
parameters. Each record supplies a centered momentum $x$, a positive
regulator $\mu\leq1$, and the same complex response coefficient
$\mathcal K(x)$. The centered-box condition $|x_j|\leq\pi$ gives
$|x|\leq2\pi$.

For the all-centered estimate, use scale $T_1=|x|+\mu$. The whole-frame
bound from the two channels is
\[
 |D-d|\leq A_{\rm int}T_1^{-4},\qquad
 |Y-c|\leq A_{\rm ref}T_1^{-4}.
\]
The exact first identity in (AT5a) then gives
$|\mathcal K|\leq(A_{\rm int}/4+A_{\rm ref}/2)T_1^{-4}$. This is the
first radial estimate; it applies also at $x=0$.

For the annular comparison, restrict the same records to
$\mu\leq|x|$ and use the second scale $T_2=|x|$. The upper channel is split
into a principal term and a tadpole term; their charged bounds are
$uT_2^{-2}$ and $v\mu^2T_2^{-6}$. The same split for the reference channel
has bounds $u_{\mathrm{ref}}T_2^{-2}$ and
$v_{\mathrm{ref}}\mu^2T_2^{-6}$. The lower-channel
remainders after subtracting their respective halos are bounded by
$A_0(T_2+\mu^2/T_2^2)T_2^{-4}$ and
$B_0(T_2+\mu^2/T_2^2)T_2^{-4}$. The free quotient calculation gives
$\Pi_G=\frac12(c_0-\frac12d_0)$, so these four estimates feed exactly
into (AT5b) and yield (AT7)--(AT10).

These scale receipts arise in the source through two separate applications
of the finite-family construction: the first applies the frame estimate to
every centered record, while the second restricts to the subfamily
$\mu\leq|x|$ before applying the shell estimate. They use the same response
coefficient and common constants. The annulus proof therefore does not
interchange a limit with a volume-dependent integral. The equality of this
analytic coefficient with the finite sampled pole-test coefficient is a
separate recording step; it is not inferred from (AT1).

\paragraph{Source boundary.}
This unit proves the normalized pole integral and the algebraic conversion
from the two explicit finite-volume scale receipts to the radial bounds.
The receipts themselves are produced by the local stationary-family proof
(`Rally.register` and `Rally.shell`) and remain a source-correspondence
obligation until that finite-record construction is expanded in the next
unit. In particular, the stationary-chart and row-price units are not
claimed here to establish these upper/tadpole and lower/halo estimates.
The later `HouseWord`/yield argument must also prove the finite recording
error decays after division by the site count, with its strip chosen before
the requested blocking side. The remaining path-history, retained-cutoff,
and flux-capacity suppliers are tracked in the evidence file.
